\subsection{\texorpdfstring{$L_V$ 的丰沛性}{行列式线丛的丰沛性}}\label{src:9.6.3-subsection}
\begin{theorem}\label{src:9.6.8}
设 $C$ 是特征零代数闭域 $\Bbbk$ 上亏格 $g\geq2$ 的光滑连通射影曲线. 设 $r>0$ 与 $d$ 为整数, $\xi$ 为次数 $d$ 的线丛. 设 $V$ 是 $C$ 上秩 $r/\gcd(r,d)$, 次数 $(r(g-1)-d)/\gcd(r,d)$ 的向量丛. 则线丛 $L_V$ 在 $M_{r,\xi}^{\mathrm{ps}}$ 上丰沛.
\end{theorem}
在证明此定理之前, 先说明它如何给出 $M_{r,d}^{\mathrm{ps}}$ 射影的第二个证明.
\begin{corollary}\label{src:9.6.9}
设 $C$ 是特征零代数闭域 $\Bbbk$ 上亏格 $g\geq2$ 的光滑连通射影曲线. 设 $r>0$ 与 $d$ 为整数, $\xi$ 为次数 $d$ 的线丛. 代数叠 $\cBun_{r,d}^{\mathrm{ss}}$ 与 $\cBun_{r,\xi}^{\mathrm{ss}}$ 有射影良模空间 $M_{r,d}^{\mathrm{ss}}$ 与 $M_{r,\xi}^{\mathrm{ss}}$.
\end{corollary}
\begin{proof}
由定理~\bookref{9.3.11} 和习题~\bookref{9.3.12}, 已知 $\cBun_{r,d}^{\mathrm{ss}}$ 与 $\cBun_{r,\xi}^{\mathrm{ss}}$ 有作为固有代数空间的良模空间 $M_{r,d}^{\mathrm{ss}}$ 与 $M_{r,\xi}^{\mathrm{ss}}$. 行列式态射 $\det:\cBun_{r,d}^{\mathrm{ss}}\to\underline{\mathcal P\!ic}^d$ 下降为良模空间的态射 $M_{r,d}^{\mathrm{ss}}\to\underline{\mathrm{Pic}}^d$, 且 $[\xi]\in\mathrm{Pic}^d$ 上的纤维为 $M_{r,\xi}^{\mathrm{ss}}$. 因为 $\underline{\mathrm{Pic}}^d$ 是同构于 $C$ 的 Jacobian 的射影簇 (定理~\bookref{7.4.65}), 根据丰沛性的开性 (\bookref{B.2.10}), 只需证明 $L_V$ 在每个纤维 $M_{r,\xi}^{\mathrm{ss}}$ 上的限制丰沛. 这正是定理~\bookref{9.6.8} 的结论.
\end{proof}
为证明定理~\bookref{9.6.8}, 由于已经知道 $L_V$ 半丰沛 (推论~\bookref{9.6.7}), 只需证明对 $m\gg0$, 态射 $|L_V^{\otimes m}|:M_{r,\xi}^{\mathrm{ps}}\to\mathbb P^N$ 拟有限. 我们将由以下结果推出这一点:
\begin{proposition}\label{src:9.6.10}
设 $E$ 与 $E'$ 是秩 $r$, 次数 $d$ 的多重稳定向量丛, 且 $E$ 有一个稳定直和项不是 $E'$ 的直和项. 存在整数 $m>0$ 及秩 $mr'$, 次数 $mr'(g-1)-md'$ 的稳定向量丛 $V$, 使
\[
\mathrm H^0(C,V\otimes E)\neq0,\qquad\mathrm H^0(C,V\otimes E')=0.
\]
特别地, $M_{r,\xi}^{\mathrm{ss}}$ 上行列式线丛 $L_V$ 的截面 $s_V$ 将至少有一个稳定直和项不同的多重稳定丛所确定的点分离.
\end{proposition}
\begin{remark}\label{src:9.6.11}
并非对任意两个不同构的多重稳定丛 $E$ 与 $E'$, 都存在适当秩与次数的向量丛 $V$, 使得 $\mathrm H^0(C,V\otimes E)\neq0$, $\mathrm H^0(C,V\otimes E')=0$. 例如 $E=E_1^{\oplus2}\oplus E_2$, $E'=E_1\oplus E_2^{\oplus2}$, 其中 $E_1$ 与 $E_2$ 是秩与次数相同的不同构稳定向量丛. 然而可以证明, 对 $m\gg0$, 形如 $s_V$ 的截面张成 $L_V$ 的整体截面向量空间, 并诱导闭嵌入 $M_{r,d}^{\mathrm{ss}}\hookrightarrow\mathbb P^N$; 见 \cite[推论 7.15]{CK07}.
\end{remark}
\begin{proof}[命题~\bookref{9.6.10} 的证明]
设 $E_1,\ldots,E_n$ 为 $E'$ 的稳定直和项, $E_0$ 为 $E$ 的一个与任何 $E_i$ 都不同构的稳定直和项. 令 $r_i:=\mathrm{rk}\,E_i$. 只需找一个秩 $mr'$, 次数 $mr'(g-1)-md'$ 的向量丛 $V$, 使
\[
\mathrm H^0(C,V\otimes E_0)\neq0,\qquad\mathrm H^0(C,V\otimes E_i)=0\quad(i=1,\ldots,n).
\]
设 $W$ 是秩 $mr'$, 次数 $mr'(g-1)-md'+1$ 的稳定向量丛, 满足下面引理~\bookref{9.6.12} 的结论 (1)--(3). 根据~\bookref{9.6.12}(1), 有非零映射 $\phi:E_0^\vee\to W$. 只需给出子层 $V\subseteq W$, 满足 $\mathrm{rk}\,V=\mathrm{rk}\,W$, $\deg V=\deg W-1$, 使 $\phi:E_0^\vee\to W$ 经过 $V$ 分解, 而对 $i=1,\ldots,n$, 无非零映射 $E_i^\vee\to W$ 经过 $V$ 分解.

根据~\bookref{9.6.12}(2), $\phi^\vee:W^\vee\to E_0$ 满射, 所以余核 $Q=\operatorname{coker}(E_0^\vee\to W)$ 局部自由. 对 $i=1,\ldots,n$ 及非零映射 $\psi_i:E_i^\vee\to W$, 诱导映射 $E_0^\vee\oplus E_i^\vee\to W$ 单射且有局部自由余核, 因为其对偶由~\bookref{9.6.12}(3) 满射. 包含关系
\[
E_0^\vee\hookrightarrow E_0^\vee\oplus E_i^\vee\xhookrightarrow{(\phi,\psi_i)}W
\]
诱导短正合列
\[
0\longrightarrow E_i^\vee\xrightarrow{\overline\psi_i}Q\longrightarrow W/(E_0^\vee\oplus E_i^\vee)\longrightarrow0.
\]
所以 $\overline\psi_i:E_i^\vee\to Q$ 单射且有局部自由余核, 从而对任何 $p\in C(\Bbbk)$, 纤维上的映射 $\overline\psi_i\otimes\kappa(p)$ 也单射.

对 $p\in C(\Bbbk)$, 考虑参数化 $Q\otimes\kappa(p)$ 中 $r_i$ 维子空间的 Grassmann 簇 $\Gr_i:=\Gr(\mathrm{rk}\,Q-r_i,Q\otimes\kappa(p))$. 记 $M_i\subset\Gr_i$ 为将非零映射 $\psi_i:E_i^\vee\to W$ 映到 $\im(\psi_i\otimes\kappa(p))$ 的映射 $\mathbb P(\mathrm{Hom}_{\mathcal O_C}(E_i^\vee,W)^\vee)\to\Gr_i$ 的像. 根据~\bookref{9.6.12}(1), $\dim M_i\leq r_i-1$.

对超平面 $H\subseteq Q\otimes\kappa(p)$, 参数化包含于 $H$ 的子空间的 Schubert 簇 $Z_{i,H}\subset\Gr_i$ 的余维为 $r_i$, 故交 $M_i\cap Z_{i,H}$ 的预期维数为 $-1$. 因此, 根据 Bertini--Kleiman 定理 \cite[推论 4(i)]{Kle74}, 对一般超平面 $H$, 当 $i=1,\ldots,n$ 时, 交 $M_i\cap Z_{i,H}$ 为空. 固定任一这样的 $H$, 考虑
\[
V:=\ker\big(W\to W\otimes\kappa(p)\to Q\otimes\kappa(p)\to(Q\otimes\kappa(p))/H\big).
\]
(这是初等修改的例子, 这一概念也出现在 Langton 定理 (\bookref{9.3.1}) 的证明中.) 按构造, $\phi:E_0^\vee\to W$ 经过 $V$ 分解, 但无非零映射 $E_i^\vee\to W$ 经过 $V$ 分解, 因为复合
\[
E_i^\vee\otimes\kappa(p)\xrightarrow{\psi_i\otimes\kappa(p)}W\otimes\kappa(p)\longrightarrow Q\otimes\kappa(p)/H
\]
非零.
\end{proof}
\begin{lemma}\label{src:9.6.12}
设 $r>0$ 与 $d$ 为整数, 置 $r'=r/\gcd(r,d)$, $d'=d/\gcd(r,d)$. 设 $E_0,E_1,\ldots,E_n$ 是固定斜率 $\mu(E_i)=d/r$, 秩 $\mathrm{rk}(E_i)=r_i$ 的稳定向量丛. 对 $m\gg0$, 秩 $mr'$, 次数 $mr'(g-1)-md'+1$ 的一般稳定向量丛 $W$ 满足以下条件:
\begin{enumerate}[(1)]
\item $\dim\mathrm{Hom}_{\mathcal O_C}(W^\vee,E_i)=r_i$ 且 $\mathrm{Ext}^1_{\mathcal O_C}(W^\vee,E_i)=0$;
\item 对 $0\leq i\leq n$, 每个非零映射 $W^\vee\to E_i$ 都满射;
\item 对 $1\leq i\leq n$ 及非零映射 $W^\vee\to E_0$ 与 $W^\vee\to E_i$, 诱导映射 $W^\vee\to E_0\oplus E_i$ 满射.
\end{enumerate}
\end{lemma}
\begin{proof}
有
\[
\dim\mathrm{Hom}_{\mathcal O_C}(W^\vee,E_i)=\mathrm h^0(W\otimes E_i)=\chi(C,W\otimes E_i)+\mathrm h^1(C,W\otimes E_i),
\]
并由 Riemann--Roch,
\[
\begin{aligned}
\chi(C,W\otimes E_i)&=\deg(W\otimes E_i)+\mathrm{rk}(W\otimes E_i)(1-g)\\
&=\deg(W)\mathrm{rk}(E_i)+\mathrm{rk}(W)\deg(E_i)+\mathrm{rk}(W)\mathrm{rk}(E_i)(1-g)\\
&=(mr'(g-1)-md'+1)r_i+mr'r_i\mu+mr'r_i(1-g)\\
&=r_i.
\end{aligned}
\]
为得到 (1), 只需证明对 $0\leq i\leq n$, $\mathrm H^1(C,W\otimes E_i)=0$. 由 Serre 对偶,
\[
\mathrm h^1(C,W\otimes E_i)=\mathrm h^0(C,W^\vee\otimes\omega_C\otimes E_i^\vee)=\dim\mathrm{Hom}_{\mathcal O_C}(W\otimes\omega_C^\vee,E_i^\vee).
\]
因为 $E_i^\vee$ 是斜率 $-d'/r'$ 的稳定丛, 而 $W^\vee\otimes\omega_C$ 是秩 $mr'$, 次数 $mr'(g-1)+md'-1$ 的一般稳定丛, 可对每个 $E_i^\vee$ 取 $\epsilon=-1$ 应用引理~\bookref{9.6.1}(2), 得对 $m\gg0$, 每个 $i$ 均有 $\mathrm{Hom}_{\mathcal O_C}(W\otimes\omega_C^\vee,E_i^\vee)=0$.

对 (2), 对每个 $E_i$ 取 $\epsilon=1$ 应用引理~\bookref{9.6.1}(1), 得到对 $m\gg0$ 及一般稳定丛 $W$, 每个非零映射 $W^\vee\to E_i$ 的像的斜率至少为 $\mu(E_i)$. 因为 $E_i$ 稳定, 每个这样的映射均满射.

对 (3), 对 $1\leq i\leq n$ 的每个 $E_0\oplus E_i$, 再取 $\epsilon=1$ 应用引理~\bookref{9.6.1}(1), 得到对 $m\gg0$ 及一般稳定丛 $W$, 每个非零映射 $W^\vee\to E_0\oplus E_i$ 的像的斜率为 $d/r$. 对任何这样的 $W$ 及两个非零映射 $W^\vee\to E_0$, $W^\vee\to E_i$, 映射 $W^\vee\to E_0\oplus E_i$ 的像斜率为 $d/r$, 且满射到两个直和项. 但 $E_0\oplus E_i$ 的斜率为 $\mu$ 的非零子层只有 $E_0$, $E_i$ 与 $E_0\oplus E_i$. 因为 $E_0$ 与 $E_i$ 不同构, 像必为整个 $E_0\oplus E_i$.
\end{proof}
\begin{proof}[定理~\bookref{9.6.8} 的证明]
根据推论~\bookref{9.6.7}, $M_{r,\xi}^{\mathrm{ps}}$ 上的行列式线丛 $L_V$ 半丰沛. 对 $m\gg0$, $L_V^{\otimes m}$ 无基点, 并定义态射 $M_{r,\xi}^{\mathrm{ps}}\to\mathbb P^N$, 其中 $N$ 为某个整数. 我们断言此映射在 $m\gg0$ 时拟有限. 事实上, 对给定的多重稳定丛 $E=E_1^{\oplus r_1}\oplus\cdots\oplus E_s^{\oplus r_s}$, 其中 $E_i$ 为两两不同构的稳定丛, 稳定直和项恰好为 $E_1,\ldots,E_s$ 的其他多重稳定丛 $E'$ 只有有限个. 对任何其他多重稳定丛 $E'$, 命题~\bookref{9.6.10} 给出秩 $m\,\mathrm{rk}\,V$, 次数 $m\deg V$ 的向量丛 $V_m$, 使 $\mathrm H^0(C,V\otimes E)\neq0$, $\mathrm H^0(C,V\otimes E')=0$. 由引理~\bookref{9.4.22}, 这定义将 $E$ 与 $E'$ 分离的截面 $s_{V_m}\in\mathrm H^0(M_{r,\xi}^{\mathrm{ps}},L_V^{\otimes m})$.

因为 $M_{r,\xi}^{\mathrm{ps}}$ 固有, 拟有限态射 $M_{r,\xi}^{\mathrm{ps}}\to\mathbb P^N$ 必有限 (推论~\bookref{5.5.14}). 这给出 $L_V^{\otimes m}$ 在 $m\gg0$ 时的丰沛性.
\end{proof}
\begin{remark}\label{src:9.6.13}
Faltings \cite[定理 I.4]{Fal93} (另见 \cite[引理 4.2]{Ses93}) 用另一种方法证明 $L_V^{\otimes m}$ 在 $m\gg0$ 时定义拟有限映射 $M_{r,\xi}^{\mathrm{ps}}\to\mathbb P^N$. 为证明此映射不压缩任何曲线, Faltings 考虑来自光滑连通射影曲线 $T$ 的映射 $T\to\cBun_{r,\xi}^{\mathrm{ss}}$, 它对应于 $C\times T$ 上的半稳定向量丛族 $\mathcal E$. 他随后证明 $\mathcal L_V|_T$ 的次数非负, 且等于零当且仅当对所有 $t\in T(\Bbbk)$, 纤维 $\mathcal E_t=\mathcal E|_{C\times\kappa(t)}$ 均 S-等价. 有趣的是, Faltings 使用了曲线 $T$ 上向量丛的稳定性, 以及将 $T$ 上 Harder--Narasimhan 滤过与 $C$ 上 Jordan--H\"older 滤过比较的跷跷板论证.
\end{remark}
